\section{結論}
\label{sec:conclusion}

本論文は、VED$\times$IFGT枠組みの最初の応用層として、概念地形の場理論的記述であるConceptual Topography（CT）を展開した。

中心となる提案は、再帰的なトークンと場の相互作用のもとで履歴によって形成された情報場$I(x,\tau)$の幾何を、モデル内の概念地形として用いることである。
この呼称は、収束、認知機構、経験へのアクセスについての主張とは別である。
モデル内では、意味の原始概念、志向的状態、意味を担う単位は必要としない。

保持する定義と提案する解釈は以下のとおりである。
\begin{enumerate}
\item 概念は、摂動に対する安定性を伴う$\nabla\Phi\approx 0$によって定義される、情報ポテンシャル$\Phi = K*I$の局所的安定領域である。
その深さ、幅、頑健性は、相互作用の履歴とともに連続的に発展する動的量である。
\item 理解は概念領域の通過コストの低下であり、認識論的な達成ではない。
地形の発展に伴って、そこに入り、そこから出ることができる。
\item 言語トークンは、意味内容を担わずに$\nabla\Phi$へ結合する、離散的で再アドレス可能な単位として機能する。
低コストな再帰がこの構成上の役割を動機づける。
反復的な相互作用による深化は定性的仮説に留まり、結合の保証された帰結ではない。
\item モデル内では、厳密な正値性$I(x,\tau)>0$を、通常の完全な概念的閉包を排するものと解釈する。
局所的な静穏は許容する。
新たな摂動やより広い観測窓のもとでのさらなる分化は、条件付きの提案であり、正値性の各点での帰結ではない。
\item 社会的スケールでは、異なるコストのもとで比較可能なトークン機能が、共有地形仮説を動機づける。
規範、制度、パラダイムをアトラクタと比較する。
深さと安定性の関係は、独立した社会的基準を必要とする仮説に留まる。
集合的な維持は局所的な機能語彙であり、確立された知性の一般的定義ではない。
\end{enumerate}

本枠組みは関連ノートシリーズを改訂したり置き換えたりしない。
記録された直観とのCT内の教育的比較を保持するが、独立に同定されたノートを検証する支持とはしない。
ノートとMaster表の正確な出典・版の帰属は未解決である。

定量的なパラメータ同定、概念個体化のトポロジー的基準、マルチエージェントの同期ダイナミクス、神経的対応といった詳細な展開は、今後の課題である。
長時間対象、結合のシグネチャ、正値性の保存、その他の形式的OPEN問題は、文章の範囲を限定したことによって解決されたわけではない。
同様に、理解を感じられる容易さから、洞察を主観的突然性から分けたことは、いずれの対応も確立しない。

要するに、Conceptual Topographyは、VED$\times$IFGT枠組みに基づき、概念構造と認知機能の場理論的表現を提案する。
概念と物理的粒子は、同一の実体や還元可能な対象としてではなく、投影された安定化領域として比較しうる。
反復する生成文法は、この比較における発見的なモチーフであり、共通方程式の定理ではない。
深さ、幅、持続、観測窓は比較を整理するが、共通量、共通場、型を明示した翻訳を確立しない。
経験的検証や主観的経験の説明は主張しない。
